\chapter{测量与时间}\label{ch:g1:measure}

哪支铅笔更长？今天是星期几？一颗糖多少钱？
本章比较长度，并用厘米来测量，
把一周的日子排好顺序，认读钟面上的整时，
并数一数小面额的硬币。

\section{长与短}

\begin{method}[比较长度]\label{met:g1:measure:compare}
比较两个长度时，让它们的\emph{一端对齐}：
\begin{enumerate}
  \item 把两个物体并排放好，从同一条线开始；
  \item 伸得更远的那个是\emph{更长}的，另一个是
        \emph{更短}的；
  \item 如果两者伸到同一处，它们的\emph{
        长度相同}。
\end{enumerate}
起点对齐是关键——从起点歪斜处比较，
很容易让眼睛看错。
\end{method}

\begin{omfigure}
\begin{tikzpicture}[scale=0.9]
  \draw[very thick, omDef] (0,0.6) -- (3.4,0.6);
  \draw[very thick, omThm] (0,0) -- (2.2,0);
  \draw[gray, dashed] (0,-0.2) -- (0,0.8);
  \node[right, font=\small] at (3.6,0.6) {更长};
  \node[right, font=\small] at (3.6,0) {更短};
\end{tikzpicture}

{\small 两根小棒左端对齐：蓝色的那根伸得更远，
所以它更长。}
\end{omfigure}

\begin{definition}[厘米]\label{def:g1:measure:cm}
要说明\emph{有多长}，用尺子以
\emph{厘米}（cm）为单位来量。把尺子的 $0$ 刻度对准物体的起点，
读出物体另一端所对的数。
\end{definition}

\begin{omfigure}
\begin{tikzpicture}[scale=1.0]
  % ruler body extends a little past the last mark so no numeral is cut
  \draw[thick] (-0.4,0) rectangle (8.4,0.7);
  \foreach \i in {0,...,8} {
    \draw[thick] (\i,0.7) -- (\i,0.4);
    \node[font=\footnotesize] at (\i,0.2) {\i};
  }
  \draw[very thick, omThm] (0,1.15) -- (5,1.15);
  \node[omThm, above, font=\small] at (2.5,1.22) {铅笔：$5$ cm};
\end{tikzpicture}

{\small 测量：$0$ 对准起点，读终点。这支铅笔
长 $5$ cm。（毫米和米见
\cref{ch:g3:measure}。）}
\end{omfigure}

\section{一周的日子}

\begin{example}[七天]\label{ex:g1:measure:days}
一周有 $7$ 天，总是按同一顺序排列：
\[
\text{星期一、星期二、星期三、星期四、星期五、星期六、
星期日。}
\]
过了星期日，又从星期一开始。\emph{昨天}是
前一天，\emph{明天}是后一天：如果今天是星期三，
那么昨天是星期二，明天是星期四。
\end{example}

\section{钟面上的整时}

\begin{method}[认读整时]\label{met:g1:measure:clock}
长针正指向 $12$ 时，这便是整时：
读出\emph{短}针所指的数。短针指向
$3$，长针指向 $12$：这是 $3$ 点。
\end{method}

\begin{omfigure}
\begin{tikzpicture}[scale=0.95]
  \draw[thick] (0,0) circle (1.3);
  \foreach \h in {1,...,12} {
    \node[font=\footnotesize] at ({90-\h*30}:1.05) {\h};
    \draw ({90-\h*30}:1.22) -- ({90-\h*30}:1.3);
  }
  \draw[very thick, omDef] (0,0) -- (0:0.62);
  \draw[thick, omThm] (0,0) -- (90:1.02);
  \node[below, font=\small] at (0,-1.55) {$3$ 点};
  \begin{scope}[xshift=4.4cm]
  \draw[thick] (0,0) circle (1.3);
  \foreach \h in {1,...,12} {
    \node[font=\footnotesize] at ({90-\h*30}:1.05) {\h};
    \draw ({90-\h*30}:1.22) -- ({90-\h*30}:1.3);
  }
  \draw[very thick, omDef] (0,0) -- ({90-8*30}:0.62);
  \draw[thick, omThm] (0,0) -- (90:1.02);
  \node[below, font=\small] at (0,-1.55) {$8$ 点};
  \end{scope}
\end{tikzpicture}

{\small 整时：长的红色指针指向 $12$，短的蓝色指针
给出时刻。（半点和一刻见
\cref{ch:g2:measure}。）}
\end{omfigure}

\section{小面额硬币}

\begin{example}[凑出一个数额]\label{ex:g1:measure:coins}
用 $1$、$2$、$5$ 的硬币，许多数额都能用几种方法凑出。
要凑出 $7$：
\[
5 + 2, \qquad
5 + 1 + 1, \qquad
2 + 2 + 2 + 1 .
\]
先取面值最大的硬币，通常是最快的方法，
所用的硬币数最少。
\end{example}

\begin{omfigure}
\begin{tikzpicture}[scale=0.85]
  \draw[thick, fill=omMeth!20] (0,0) circle (0.5);
  \node[font=\small] at (0,0) {$5$};
  \node[font=\large] at (1.0,0) {$+$};
  \draw[thick, fill=omDef!20] (2.0,0) circle (0.42);
  \node[font=\small] at (2.0,0) {$2$};
  \node[font=\large] at (3.0,0) {$=$};
  \node[font=\small] at (4.0,0) {$7$};
\end{tikzpicture}

{\small 面值 5 的硬币和面值 2 的硬币合起来是 7。}
\end{omfigure}

\section{练习}

\begin{exercise}[$\star$]\label{exo:g1:measure:1}
把三支铅笔一端对齐排成一排。哪一支最长？
哪一支最短？
\end{exercise}

\begin{exercise}[$\star$]\label{exo:g1:measure:2}
用直尺量一量：这本书的长边、你的小指、
一把勺子。请以整厘米作答。
\end{exercise}

\begin{exercise}[$\star$]\label{exo:g1:measure:3}
说出下列日期：星期一的后一天；星期日的前一天；星期五的后两天。
\end{exercise}

\begin{exercise}[$\star$]\label{exo:g1:measure:4}
今天是星期四。昨天是星期几？明天是星期几？
三天后是星期几？
\end{exercise}

\begin{exercise}[$\star$]\label{exo:g1:measure:5}
画一个显示 $6$ 点的钟面。时针在哪里？
分针在哪里？
\end{exercise}

\begin{exercise}[$\star$]\label{exo:g1:measure:6}
读出整点：时针指向 $5$，分针指向 $12$；时针指向 $11$，
分针指向 $12$。
\end{exercise}

\begin{exercise}[$\star$]\label{exo:g1:measure:7}
用两种不同的方法凑出 $8$，硬币面值为 $1$、$2$ 和 $5$。
\end{exercise}

\begin{exercise}[$\star$]\label{exo:g1:measure:8}
莱奥有一枚面值 $5$ 的硬币和三枚面值 $2$ 的硬币，
他一共有多少钱？
\end{exercise}

\begin{exercise}[$\star\star$]\label{exo:g1:measure:9}
一颗糖果售价 $3$，佐伊用一枚面值 $5$ 的硬币付款，
她应找回多少钱？（从 $3$ 数到 $5$。）
\end{exercise}
